\scanpage{154}
\mspage{76}
On \deleted{désigne donc par $D_{d,d'}^\Sigma$ l'image} \inserted{posera donc}
\[
(250)\qquad D_{d,d'}^\Sigma=\text{image inverse de }D_{d'}^\Sigma\text{ par}
\]
\[
D_{d,d'}\longrightarrow D_{d'}\quad\text{(et \uline{non} l'image inv.}
\]
\[
\text{de }D_d^\Sigma\text{ par }D_{d,d'}\longrightarrow D_d\ !\ )
\]
de sorte que \uline{$D_{d,d'}^\Sigma$} est aussi une réunion de comp. conn. de $D_{d,d'}$, et on a
\[
(251)
\]
\diagram{\begin{tikzcd}[ampersand replacement=\&,column sep=small]
\&D_{d,d'}^\Sigma\arrow[dl]\arrow[dr]\&\\
D_d^\Sigma\&\&D_{d'}^\Sigma
\end{tikzcd}}.

La donnée d'une partie $\Theta$-admissible $\Sigma$ de $X=D_0$ est équivalente à la donnée d'un système de $D_d^\Sigma$, (équivalente à la donnée de parties correspondantes
\[
\deleted{\pi_0}(252)\qquad
\pi_0(D_d^\Sigma)\overset{\mathrm{def}}=\pi_0(D_d)^\Sigma\subset\pi_0(D_d)),
\]
satisfaisant les conditions de compatibilité (251).

On désignera alors, pour tout drapeau $d_*=\{d_0>d_1>\cdots>d_r\geq0\}$, par
\[
(253)\qquad D_{d_0\ldots d_r}^\Sigma=\text{image inverse de }D_{d_r}^\Sigma
\text{ dans }D_{d_0,\ldots,d_r}
\]
et on voit que \inserted{$D_{\Delta_r}^\Sigma=\coprod$} $D_{d_0\ldots d_r}^\Sigma$ forme un système $\frac12$ simplicial de \unsure{telles} multiplicités analytiques, et qui, \deleted{\unsure{est un}} correspond d'ailleurs à une sous-catégorie $\mathcal D_\bullet^\Sigma$ de la
\margin{153}
